% !TEX program = lualatex % (can also use "xelatex") 
\documentclass[12pt]{article}
\input{ejoi_structure.tex}

\usepackage[bulgarian]{babel}

\begin{document}

\renewcommand{\round}{Ден 2}
\renewcommand{\problemName}{Задача Грейдър67}
\renewcommand{\lang}{Български}

\renewcommand{\tl}{$0.1$ sec.}
\renewcommand{\ml}{$256$ MB}
\problem{\problemName}
(Поради технически проблеми задачата не може да се казва grader.)\\
Дадено е число $N$. Написана е програма, която връща редица $a$ с $N$ елемента, който изпълнява следните условия:
\begin{itemize}
    \item има поне една двойка индекси $1 \leq i,j \leq N, i \neq j$, където $a_i \neq a_j$.
	\item $\sum_{i=1}^{N} a_i \equiv 0 \pmod{N}$.
	\item разликата $max(a_1, a_2, ...) - min(a_1, a_2, ...)$ е минимална.
\end{itemize}
\subsection{Детайли по имплементацията}
Вашата задача е да напишете грейдър, който проверява коректността на даденото решение. То е имплементирано чрез функцията:
\begin{lstlisting}
 std::vector<int> solve(int N);
\end{lstlisting}
\textbf{Не трябва} да я имплементирате! Трябва да имплементирате функцията \texttt{main}, която въвежда $N$, извиква \texttt{solve} и проверява дали функцията връща правилно решение. Не е гарантирано, че върнатата редица отговаря на всички условия. Грейдърът трябва да отпечатва едно число между 0 и 1 - точките на даден тест. Точките се смятат по следния начин:
\begin{itemize}
    \item изпълнено първо условие носи 0.15 на тест;
	\item изпълнено второ условие носи 0.25 на тест;
	\item изпълнено трето условие носи 0.60 на тест.
\end{itemize}

\subsection{Ограничения}
\begin{itemize}
\item $1 \leq N \leq 10^5+67$
\end{itemize}

\subsection{Подзадачи}
\begin{table}[H]
\begin{tblr}{|Q[c,m]|Q[c,m]|Q[c,m]|}
\hline
\textbf{Подзадача} & \textbf{Точки} & \textbf{Ограничение}\\
\hline
$1$ & $15$ & $1 \leq N \leq $ около $6-7$\\
\hline
$2$ & $20$ & $1 \leq N \leq 30$\\
\hline
$3$ & $25$ & $1 \leq N \leq 100$\\
\hline
$4$ & $20$ & $1 \leq N \leq 1067$\\
\hline
$5$ & $20$ & $1 \leq N \leq 10^5+67$\\
\hline
\end{tblr}
\end{table}
\FloatBarrier

\subsection{Пример}
\begin{table}[H]
\begin{tblr}{|Q[c,m]|Q[c,m]|}
\hline
\textbf{Жури} & \textbf{Състезател}\\
\hline
$ $ & \texttt{solve(5)}\\
\hline
${2, 1, 3, 2, 3}$ & Грейдърът връща 0.75.\\
\hline
\end{tblr}
\end{table}
Този грейдър получава пълни точки на теста, защото правилно открива точките, които получава върнатата редица.
\FloatBarrier

\subsection{Примерен грейдър}
Предоставен Ви е файл \texttt{Lgrader.cpp}, който действа като грейдър на Вашия грейдър. За да използвате функцията за комуникацията, във Вашия код трябва да включите библиотеката \texttt{grader67.h}:
\begin{lstlisting}
 #include "grader67.h"
\end{lstlisting}
Форматът на входа е следния: 
\begin{itemize}
\item ред $1$: числото $N$.
\item ред 2: масива $a$, върху който ще се грейдва Вашия грейдър. 
\end{itemize}

След това се изпълнява Вашия грейдър, който извежда точките, които получава това решение.

\end{document}

